\clearpage
\begin{proofbox}
\textbf{证明前先分清坐标：旧点叫 $(u,v)$，新点叫 $(X,Y)$。}
旧图上有 $v=f(u)$。先按几何对称找到 $X,Y$，再消去 $u,v$，最后把新坐标 $X,Y$ 改写为习惯的 $x,y$。不能把变换前后的横、纵坐标同时叫 $x,y$，却又当成没有改变的数。
\medskip
\textbf{步骤一（横轴对称）}
\begin{center}\begin{tikzpicture}[font=\small,>=stealth,line join=round]
\draw[gray!40,rounded corners=3pt] (0,0) rectangle (14.2,8.3);
\node[text width=12.6cm,align=center] at (7.1,7.65) {关于 $x$ 轴：同一横坐标，上下高度相反};
\begin{scope}[shift={(4.3,4.1)},x=1.65cm,y=0.47000000000000003cm]
\draw[->,gray] (-2,0)--(2,0) node[right]{$x$};\draw[->,gray] (0,-5)--(0,5) node[above]{$y$};

\begin{scope}\clip (-2,-5) rectangle (2,5);
\draw[blue,thick,domain=-2:2,samples=180] plot(\x,{3*\x+2});
\draw[orange!85!black,thick,domain=-2:2,samples=180] plot(\x,{-3*\x-2});
\end{scope}
\fill[blue] (0,2) circle(2.2pt);\node[above left=6pt,blue] at (0,2) {$P$};
\fill[orange!85!black] (0,-2) circle(2.2pt);\node[below left=6pt,orange!85!black] at (0,-2) {$Q$};
\end{scope}
\node[anchor=west,align=left,text width=4.8cm] at (8.65,4.5) {蓝：$f(x)=3x+2$\\[8pt]橙：$g(x)=-3x-2$\\[8pt]$P(0,2)\to Q(0,-2)$\\[8pt]$g(x)=-f(x)$};
\node[text width=12.6cm,align=center] at (7.1,.65) {两幅图关于横轴相互对称；原函数本身不必关于横轴对称。};
\end{tikzpicture}\end{center}


旧点 $(u,v)$ 满足 $v=f(u)$。关于 $x$ 轴反射后，$X=u,Y=-v$，故
\[Y=-v=-f(u)=-f(X).\]
于是新函数为 $g(x)=-f(x)$，定义域仍为 $D$。反过来，新图上 $(X,-f(X))$ 反射后回到 $(X,f(X))$，所以对应覆盖两幅完整图像。

\end{proofbox}
\clearpage
\begin{proofbox}
\textbf{步骤二（纵轴对称）}
\begin{center}\begin{tikzpicture}[font=\small,>=stealth,line join=round]
\draw[gray!40,rounded corners=3pt] (0,0) rectangle (14.2,8.3);
\node[text width=12.6cm,align=center] at (7.1,7.65) {关于 $y$ 轴：左右位置相反，高度相等};
\begin{scope}[shift={(4.3,2.1772727272727272)},x=0.7333333333333333cm,y=0.4272727272727273cm]
\draw[->,gray] (-4.5,0)--(4.5,0) node[right]{$x$};\draw[->,gray] (0,-1)--(0,10) node[above]{$y$};
\draw[violet,dashed] (-2,1)--(2,1);
\begin{scope}\clip (-4.5,-1) rectangle (4.5,10);
\draw[blue,thick,domain=-4.5:4.5,samples=180] plot(\x,{\x*\x-4*\x+5});
\draw[orange!85!black,thick,domain=-4.5:4.5,samples=180] plot(\x,{\x*\x+4*\x+5});
\end{scope}
\fill[blue] (2,1) circle(2.2pt);\node[above right=5pt,blue] at (2,1) {$P$};
\fill[orange!85!black] (-2,1) circle(2.2pt);\node[above left=5pt,orange!85!black] at (-2,1) {$Q$};
\end{scope}
\node[anchor=west,align=left,text width=4.8cm] at (8.65,4.5) {蓝：$f(x)=(x-2)^2+1$\\[8pt]橙：$g(x)=(x+2)^2+1$\\[8pt]$P(2,1)\to Q(-2,1)$\\[8pt]$g(x)=f(-x)$};
\node[text width=12.6cm,align=center] at (7.1,.65) {只显示窗口内曲线；每个原图上的点 $(u,v)$ 都对应另图上的 $(-u,v)$。};
\end{tikzpicture}\end{center}


旧点 $(u,v)$ 反射后，$X=-u,Y=v$。因此 $u=-X$，代入 $v=f(u)$：
\[Y=f(-X).\]
新函数为 $g(x)=f(-x)$，定义域是 $\{-u:u\in D\}$。反射两次又回到原来的点，故新图中的每一点也都来自旧图。

\end{proofbox}
\clearpage
\begin{proofbox}
\textbf{步骤三（原点中心对称）}
\begin{center}\begin{tikzpicture}[font=\small,>=stealth,line join=round]
\draw[gray!40,rounded corners=3pt] (0,0) rectangle (14.2,8.3);
\node[text width=12.6cm,align=center] at (7.1,7.65) {关于原点：横、纵坐标同时取相反数};
\begin{scope}[shift={(4.3,4.1)},x=1.65cm,y=0.5875cm]
\draw[->,gray] (-2,0)--(2,0) node[right]{$x$};\draw[->,gray] (0,-4)--(0,4) node[above]{$y$};
\draw[violet,dashed] (-1,-1)--(1,1);
\begin{scope}\clip (-2,-4) rectangle (2,4);
\draw[blue,thick,domain=-2:2,samples=180] plot(\x,{2*\x-1});
\draw[orange!85!black,thick,domain=-2:2,samples=180] plot(\x,{2*\x+1});
\end{scope}
\fill[blue] (1,1) circle(2.2pt);\node[above left=6pt,blue] at (1,1) {$P$};
\fill[orange!85!black] (-1,-1) circle(2.2pt);\node[below right=6pt,orange!85!black] at (-1,-1) {$Q$};
\end{scope}
\node[anchor=west,align=left,text width=4.8cm] at (8.65,4.5) {蓝：$f(x)=2x-1$\\[8pt]橙：$g(x)=2x+1$\\[8pt]$P(1,1)\to Q(-1,-1)$\\[8pt]$g(x)=-f(-x)$};
\node[text width=12.6cm,align=center] at (7.1,.65) {连线的中点是原点；两条直线相互对称，并不要求它们重合。};
\end{tikzpicture}\end{center}


旧点 $(u,v)$ 关于原点的对称点为 $(X,Y)=(-u,-v)$。从 $u=-X,v=-Y$ 和 $v=f(u)$ 得
\[-Y=f(-X),\qquad Y=-f(-X).\]
所以 $g(x)=-f(-x)$，定义域为 $\{-u:u\in D\}$；反向做同一对称也回到旧图。

\end{proofbox}
\clearpage
\begin{proofbox}
\textbf{步骤四（反函数对称）}
\begin{center}\begin{tikzpicture}[font=\small,>=stealth,line join=round]
\draw[gray!40,rounded corners=3pt] (0,0) rectangle (14.2,8.3);
\node[text width=12.6cm,align=center] at (7.1,7.65) {关于 $y=x$：交换两坐标，得到反函数图像};
\begin{scope}[shift={(1.94,2.69)},x=0.9400000000000001cm,y=0.9400000000000001cm]
\draw[->,gray] (-1,0)--(4,0) node[right]{$x$};\draw[->,gray] (0,-1)--(0,4) node[above]{$y$};
\draw[violet,dashed] (-.8,-.8)--(3.8,3.8) node[above]{$y=x$};\draw[gray,dashed] (2,3)--(3,2);
\begin{scope}\clip (-1,-1) rectangle (4,4);
\draw[blue,thick,domain=-1:4,samples=180] plot(\x,{2*\x-1});
\draw[orange!85!black,thick,domain=-1:4,samples=180] plot(\x,{(\x+1)/2});
\end{scope}
\fill[blue] (2,3) circle(2.2pt);\node[above left=5pt,blue] at (2,3) {$P$};
\fill[orange!85!black] (3,2) circle(2.2pt);\node[below right=5pt,orange!85!black] at (3,2) {$Q$};
\end{scope}
\node[anchor=west,align=left,text width=4.8cm] at (8.65,4.5) {蓝：$f(x)=2x-1$\\[8pt]橙：$f^{-1}(x)=\frac{x+1}{2}$\\[8pt]$P(2,3)\to Q(3,2)$\\[8pt]反函数的定义域\\是原函数的值域。};
\node[text width=12.6cm,align=center] at (7.1,.65) {横纵单位长度相同；任意一对交换坐标的点关于 $y=x$ 对称。};
\end{tikzpicture}\end{center}


旧点 $(u,v)$ 反射为 $(X,Y)=(v,u)$，所以
\[X=f(Y).\]
若 $f$ 在 $D$ 上不同输入有不同输出，那么每个 $X\in f(D)$ 恰好对应一个 $Y\in D$，可记作 $Y=f^{-1}(X)$。故新图为 $y=f^{-1}(x)$，定义域是 $f(D)$。反射两次复原，覆盖双方全部点。

\end{proofbox}
\clearpage
\begin{proofbox}
\textbf{推广证明一｜竖直轴 $x=a$}
\begin{center}\begin{tikzpicture}[font=\small,>=stealth,line join=round]
\draw[gray!40,rounded corners=3pt] (0,0) rectangle (14.2,8.3);
\node[text width=12.6cm,align=center] at (7.1,7.65) {一般竖直轴 $x=2$：横坐标之和为 $4$，高度相等};
\begin{scope}[shift={(2.0999999999999996,2.22)},x=1.0999999999999999cm,y=0.47000000000000003cm]
\draw[->,gray] (-1,0)--(5,0) node[right]{$x$};\draw[->,gray] (0,-1)--(0,9) node[above]{$y$};
\draw[violet,dashed] (2,-.7)--(2,8.5) node[above]{$x=2$};\draw[gray,dashed] (1,1)--(3,1);
\begin{scope}\clip (-1,-1) rectangle (5,9);
\draw[blue,thick,domain=-1:5,samples=180] plot(\x,{(\x-1)^2+1});
\draw[orange!85!black,thick,domain=-1:5,samples=180] plot(\x,{(\x-3)^2+1});
\end{scope}
\fill[blue] (1,1) circle(2.2pt);\node[above left=5pt,blue] at (1,1) {$P$};
\fill[orange!85!black] (3,1) circle(2.2pt);\node[above right=5pt,orange!85!black] at (3,1) {$Q$};
\end{scope}
\node[anchor=west,align=left,text width=4.8cm] at (8.65,4.5) {蓝：$f(x)=(x-1)^2+1$\\[8pt]橙：$g(x)=(x-3)^2+1$\\[8pt]$P(1,1)\to Q(3,1)$\\[8pt]$g(x)=f(4-x)$};
\node[text width=12.6cm,align=center] at (7.1,.65) {一般形式：关于 $x=a$ 相互对称的另一个函数为 $g(x)=f(2a-x)$。};
\end{tikzpicture}\end{center}


关于竖直轴反射，横坐标平均为 $a$，纵坐标相等：
\[X=2a-u,\quad Y=v.\]
由 $u=2a-X$ 和 $v=f(u)$ 得 $Y=f(2a-X)$，所以 $g(x)=f(2a-x)$。
所有旧横坐标 $u\in D$ 反射后组成 $E=\{2a-u:u\in D\}$；再反射一次回到 $u$，故充分性与必要性均成立。

\end{proofbox}
\clearpage
\begin{proofbox}
\textbf{推广证明二｜水平轴 $y=b$}
\begin{center}\begin{tikzpicture}[font=\small,>=stealth,line join=round]
\draw[gray!40,rounded corners=3pt] (0,0) rectangle (14.2,8.3);
\node[text width=12.6cm,align=center] at (7.1,7.65) {一般水平轴 $y=3$：横坐标相同，高度之和为 $6$};
\begin{scope}[shift={(1.66,2.69)},x=1.3199999999999998cm,y=0.47000000000000003cm]
\draw[->,gray] (-0.5,0)--(4.5,0) node[right]{$x$};\draw[->,gray] (0,-2)--(0,8) node[above]{$y$};
\draw[violet,dashed] (-.3,3)--(4.2,3) node[right]{$y=3$};\draw[gray,dashed] (2,1)--(2,5);
\begin{scope}\clip (-0.5,-2) rectangle (4.5,8);
\draw[blue,thick,domain=-0.5:4.5,samples=180] plot(\x,{(\x-2)^2+1});
\draw[orange!85!black,thick,domain=-0.5:4.5,samples=180] plot(\x,{5-(\x-2)^2});
\end{scope}
\fill[blue] (2,1) circle(2.2pt);\node[below right=5pt,blue] at (2,1) {$P$};
\fill[orange!85!black] (2,5) circle(2.2pt);\node[above right=5pt,orange!85!black] at (2,5) {$Q$};
\end{scope}
\node[anchor=west,align=left,text width=4.8cm] at (8.65,4.5) {蓝：$f(x)=(x-2)^2+1$\\[8pt]橙：$g(x)=5-(x-2)^2$\\[8pt]$P(2,1)\to Q(2,5)$\\[8pt]$g(x)=6-f(x)$};
\node[text width=12.6cm,align=center] at (7.1,.65) {一般形式：关于 $y=b$ 相互对称的另一个函数为 $g(x)=2b-f(x)$。};
\end{tikzpicture}\end{center}


关于水平轴反射，横坐标相同，纵坐标平均为 $b$：
\[X=u,\quad Y=2b-v.\]
代入 $v=f(u)$ 得 $Y=2b-f(X)$，所以 $g(x)=2b-f(x)$；横坐标不变，定义域仍为 $D$。再做一次反射回到原图，故是完整的双向对应。

\end{proofbox}
\clearpage
\begin{proofbox}
\textbf{推广证明三｜中心 $(a,b)$}
\begin{center}\begin{tikzpicture}[font=\small,>=stealth,line join=round]
\draw[gray!40,rounded corners=3pt] (0,0) rectangle (14.2,8.3);
\node[text width=12.6cm,align=center] at (7.1,7.65) {一般中心 $(2,3)$：连线中点固定，两坐标都取平均};
\begin{scope}[shift={(1.66,2.69)},x=1.3199999999999998cm,y=0.47000000000000003cm]
\draw[->,gray] (-0.5,0)--(4.5,0) node[right]{$x$};\draw[->,gray] (0,-2)--(0,8) node[above]{$y$};
\draw[violet,dashed] (1,1)--(3,5);
\begin{scope}\clip (-0.5,-2) rectangle (4.5,8);
\draw[blue,thick,domain=-0.5:4.5,samples=180] plot(\x,{(\x-1)^2+1});
\draw[orange!85!black,thick,domain=-0.5:4.5,samples=180] plot(\x,{5-(\x-3)^2});
\end{scope}
\fill[blue] (1,1) circle(2.2pt);\node[blue] at (1.9,.65) {$P$};
\fill[orange!85!black] (3,5) circle(2.2pt);\node[above left=6pt,orange!85!black] at (3,5) {$Q$};
\fill[violet] (2,3) circle(2.2pt);\node[below right=5pt,violet] at (2,3) {$C$};
\end{scope}
\node[anchor=west,align=left,text width=4.8cm] at (8.65,4.5) {蓝：$f(x)=(x-1)^2+1$\\[8pt]橙：$g(x)=5-(x-3)^2$\\[8pt]$P(1,1)\to Q(3,5)$\\[8pt]$g(x)=6-f(4-x)$};
\node[text width=12.6cm,align=center] at (7.1,.65) {一般形式：关于 $(a,b)$ 相互对称，$g(x)=2b-f(2a-x)$。};
\end{tikzpicture}\end{center}


中心是旧点与新点连线的中点，所以
\[X=2a-u,\quad Y=2b-v.\]
先从第一式得到 $u=2a-X$，再将 $v=f(u)$ 代入第二式：
\[Y=2b-f(2a-X).\]
故 $g(x)=2b-f(2a-x)$，定义域为 $E=\{2a-u:u\in D\}$；再转半圈复原，因此两个方向均成立。

\end{proofbox}
\clearpage
\begin{proofbox}
\textbf{图片最后一行的证明｜为什么是差的一半？}
\begin{center}\begin{tikzpicture}[font=\small,>=stealth,line join=round]
\draw[gray!40,rounded corners=3pt] (0,0) rectangle (14.2,8.3);
\node[text width=12.6cm,align=center] at (7.1,7.65) {错位输入：两个顶点的中间位置是 $2$，不是 $3$};
\begin{scope}[shift={(2.32,2.3375)},x=0.6599999999999999cm,y=0.5875cm]
\draw[->,gray] (-2,0)--(8,0) node[right]{$x$};\draw[->,gray] (0,-1)--(0,7) node[above]{$y$};
\draw[violet,dashed] (2,-.7)--(2,6.5);\node[violet] at (2,7.6){$x=2$};\draw[gray,dashed] (0,1)--(4,1);
\begin{scope}\clip (-2,-1) rectangle (8,7);
\draw[blue,thick,domain=-2:8,samples=180] plot(\x,{(1+\x)^2});
\draw[orange!85!black,thick,domain=-2:8,samples=180] plot(\x,{(5-\x)^2});
\end{scope}
\fill[blue] (-1,0) circle(2.2pt);\node[below=5pt,blue] at (-1,0) {$-1$};
\fill[orange!85!black] (5,0) circle(2.2pt);\node[below=5pt,orange!85!black] at (5,0) {$5$};
\fill[blue] (0,1) circle(2.2pt);\node[above left=5pt,blue] at (0,1) {$P$};
\fill[orange!85!black] (4,1) circle(2.2pt);\node[above right=5pt,orange!85!black] at (4,1) {$Q$};
\end{scope}
\node[anchor=west,align=left,text width=4.8cm] at (8.65,4.5) {取 $f(u)=u^2,a=1,b=5$\\[8pt]蓝：$F(x)=f(1+x)$\\橙：$G(x)=f(5-x)$\\[8pt]顶点在 $-1$ 与 $5$\\[8pt]轴：$\frac{5-1}{2}=2$};
\node[text width=12.6cm,align=center] at (7.1,.65) {一般公式是 $x=(b-a)/2$；把 $F$ 的输入位置反射过去，恰好成为 $G$ 的位置。};
\end{tikzpicture}\end{center}


令 $F(x)=f(a+x)$，$G(x)=f(b-x)$，并设 $c=(b-a)/2$。将第一幅图关于 $x=c$ 反射，得到的函数为
\[F(2c-x)=f(a+2c-x)=f(b-x)=G(x).\]
这里用到 $a+2c=b$。所以两图关于 $x=(b-a)/2$ 相互对称。

还要核对范围：第一图的横坐标集合为 $D_1=\{u-a:u\in D\}$，反射后变成
\[\{2c-(u-a):u\in D\}=\{b-u:u\in D\}=D_2,\]
恰好是第二图的定义域。因而此结论对任意原定义域 $D$ 成立，不要求 $f$ 是奇函数、偶函数或周期函数。

\end{proofbox}
